\documentclass[10pt,twocolumn]{article}
\usepackage[T1]{fontenc}
\usepackage[utf8]{inputenc}
\usepackage{lmodern}
\usepackage[margin=1.8cm,top=2.4cm,bottom=2cm,columnsep=0.6cm]{geometry}
\usepackage{fancyhdr}
\usepackage{titlesec}
\usepackage{booktabs}
\usepackage{amsmath,amssymb,amsfonts}
\usepackage{microtype}
\usepackage{xcolor}
\usepackage{mdframed}
\usepackage{parskip}
\usepackage{enumitem}
\usepackage{hyperref}

\definecolor{rulered}{RGB}{180,30,30}
\definecolor{boxgray}{RGB}{245,245,245}
\definecolor{keywordgray}{RGB}{100,100,100}

\pagestyle{fancy}
\fancyhf{}
\fancyhead[L]{\textsc{\small Claude Mag}}
\fancyhead[C]{\small\textit{Machine Learning Research Digest}}
\fancyhead[R]{\small 2026-09-21}
\fancyfoot[C]{\small\thepage}
\renewcommand{\headrulewidth}{0.8pt}

\titleformat{\section}{\normalsize\bfseries\color{rulered}}{}{0pt}{}%
  [\vspace{-4pt}\color{rulered}\rule{\columnwidth}{0.4pt}]
\titlespacing{\section}{0pt}{8pt}{4pt}

\hypersetup{colorlinks=true,urlcolor=rulered,linkcolor=black}

\begin{document}
\setlength{\parindent}{0pt}
\setlength{\parskip}{4pt}

{\small\color{keywordgray}\textbf{Keywords:}
  agentic reinforcement learning \textbullet{}
  on-policy distillation \textbullet{}
  GRPO \textbullet{}
  adaptive training schedule}\par\vspace{4pt}

{\LARGE\bfseries\raggedright
  RetireOPD}\par\vspace{4pt}

{\large\itshape\raggedright
  Adaptive Retirement of Teacher Distillation Unlocks an 18-Point Gain for
  Multi-Turn Agentic RL on ALFWorld and WebShop}\par\vspace{6pt}

{\small\textbf{By} Yan Yu, Zhengxi Lu, Yizhou Liu, Yichen Pan, Aozhe Wang,
  Qipeng Chen, Hua Yang, Wenqi Zhang, Weiming Lu, Qianglong Chen
  \& Yongliang Shen (Zhejiang University; Alibaba Group)}\par\vspace{2pt}

{\small\color{keywordgray}arXiv:2609.20784 \textbullet{} Submitted September 17, 2026}
\par\vspace{2pt}\noindent{\color{rulered}\rule{\columnwidth}{0.6pt}}\vspace{6pt}

Teacher distillation can harm as much as it helps when training agentic
reinforcement learning---unless the student knows when to stop listening.
\textit{RetireOPD} separates a skill-conditioned teacher from a skill-free
student, trains them in sequence, and lets the student autonomously retire
distillation once it has absorbed what the teacher can teach.
The result: up to $+18.8\%$ success rate on ALFWorld and $+19.0\%$ on
WebShop over pure RL across three model scales.

\begin{mdframed}[backgroundcolor=boxgray,linecolor=rulered,linewidth=1pt,
    innerleftmargin=6pt,innerrightmargin=6pt,
    innertopmargin=6pt,innerbottommargin=6pt]
\textbf{\small AT A GLANCE}\par\vspace{4pt}
\begin{description}[leftmargin=0pt,labelsep=4pt,itemsep=2pt,parsep=0pt,
    font=\small\bfseries]
  \item[Problem:] Multi-turn agents trained with sparse-reward RL receive
    weak learning signals; on-policy distillation (OPD) from a privileged
    teacher helps early but harms later when the teacher becomes a bottleneck.
  \item[Why it matters:] Web and household agents must make long chains of
    tool calls; sparse end-of-episode rewards provide almost no signal for
    individual action improvements.
  \item[Method:] Pre-train a skill-conditioned teacher with GRPO; jointly
    train a skill-free student with RL\,+\,OPD; retire OPD automatically
    once the student's discrepancy stops shrinking and it reaches a target
    success fraction of the teacher.
  \item[Result:] $+14.1$--$18.8\%$ on ALFWorld; $+11.8$--$19.0\%$ on
    WebShop vs.\ RL-only at three model scales (1.5B, 3B, 7B).
  \item[Evidence:] Ablations over teacher pre-training, retirement
    triggers, and $\alpha$ threshold sensitivity; comparisons vs.\
    SFT, GRPO, continuous OPD, and SAGE-OPD.
\end{description}
\end{mdframed}

\section{The Big Picture}

Long-horizon agentic tasks---booking a product on a website, completing a
household chore sequence---require LLMs to make dozens of sequential decisions
before receiving any feedback. Reinforcement learning can discover effective
strategies from this sparse reward but converges slowly because the credit
assignment problem is severe: which of the 40 actions in a successful
trajectory were actually critical?

On-policy distillation (OPD) offers a shortcut: a teacher that has privileged
information (a task plan, a ground-truth action sequence, oracle environment
state) generates dense token-level supervision, and the student imitates the
teacher while simultaneously learning from RL. Prior work treats OPD as a
fixed co-training objective. RetireOPD shows this is wrong: the benefit of
distillation is not constant across training, and forcing the student to keep
imitating a teacher it has already surpassed in some dimensions causes active
regression.

\section{Why Existing Approaches Fall Short}

\textbf{Pure RL (GRPO)} trains with a clipped surrogate policy gradient but
has no mechanism to bootstrap from the teacher's privileged knowledge. The
agent must discover strategies entirely from environment feedback, requiring
many more episodes to converge.

\textbf{Continuous joint OPD\,+\,RL} combines distillation and RL throughout
training. Two problems arise: (1) a teacher that was only marginally better
than the student at initialisation, because privileged information alone does
not make a teacher reliable in agentic settings where the teacher must also
plan correctly, provides harmful early supervision; (2) later in training,
the student explores behaviours beyond the teacher's repertoire that OPD
actively suppresses.

\textbf{SAGE-OPD} uses a predefined schedule to taper OPD linearly, which
is better than constant OPD but cannot adapt to the variable rate at which
different model scales or task complexities saturate distillation.

\section{Core Mechanism}

\textbf{Phase 1: Teacher pre-training.}
A skill-conditioned teacher $\pi_T(a_t \mid s_t, z)$ is initialised from
the base LLM and fine-tuned with GRPO on environment rewards augmented by
a skill-following term that rewards use of the provided privilege $z$
(e.g., sub-task labels, reference trajectories). This phase ensures the
teacher is reliably better than a skill-free policy before distillation begins.

\textbf{Phase 2: Student joint training.}
A skill-free student $\pi_S(a_t \mid s_t)$ is trained with:
\[
\mathcal{L}_S = \mathcal{L}_{\text{GRPO}}(\pi_S)
              + \beta \cdot \mathcal{L}_{\text{OPD}}(\pi_T, \pi_S),
\]
with teacher parameters frozen.

\textbf{Adaptive Retirement.}
At each training step, two signals are monitored:
\begin{itemize}[itemsep=1pt,parsep=0pt]
  \item \emph{Discrepancy:} $\Delta_t = \mathrm{KL}(\pi_T \| \pi_S)$
    averaged over recent rollouts.
  \item \emph{Success gap:} $\rho_t = R_S^t / R_T$, the student's rolling
    success rate relative to the teacher's fixed rate.
\end{itemize}
Retirement fires when:
\[
\Delta_t - \Delta_{t-k} < \epsilon_\Delta \quad \text{AND} \quad \rho_t \geq \alpha.
\]
After retirement, $\beta{=}0$ and only $\mathcal{L}_{\text{GRPO}}$ is used.

\section{Technical Formulation}

The GRPO loss is:
\[
\mathcal{L}_{\text{GRPO}} = -\mathbb{E}_\tau\sum_t \min\!\left(
  \frac{\pi_S(a_t)}{\pi_{\text{ref}}(a_t)}\hat{A}_t,\;
  \text{clip}\!\left(\frac{\pi_S}{\pi_{\text{ref}}}, 1\pm\epsilon\right)\hat{A}_t
\right)
\]
where $\hat{A}_t$ is the group-relative advantage and $\pi_{\text{ref}}$ is
the reference (SFT checkpoint or retirement-point snapshot).

The OPD loss is:
\[
\mathcal{L}_{\text{OPD}} =
  \mathbb{E}_\tau\sum_t
  \mathrm{KL}\!\left(\pi_T(\cdot \mid s_t, z)\,\big\|\,\pi_S(\cdot\mid s_t)\right).
\]

The teacher pre-training loss adds a skill-following reward:
\[
r_T = r_{\text{env}} + \lambda_z \cdot \mathbf{1}\!\left[\pi_T \text{ uses } z\right].
\]

Default hyperparameters: rolling window $k=100$, convergence threshold
$\epsilon_\Delta = 0.01$, success-fraction threshold $\alpha = 0.85$,
distillation weight $\beta = 0.5$.

\section{Learning or Inference Procedure}

\begin{enumerate}[itemsep=2pt,parsep=0pt]
  \item Pre-train skill-conditioned teacher with GRPO\,+\,skill reward for
    $T_{\text{teacher}}$ steps until teacher success rate saturates.
  \item Initialise student from base LLM (same checkpoint as teacher init).
  \item Train student jointly with $\mathcal{L}_{\text{GRPO}} + \beta\mathcal{L}_{\text{OPD}}$.
  \item At each step compute $\Delta_t$ and $\rho_t$; fire retirement
    condition if met; set $\beta{=}0$ and continue RL-only.
  \item At inference: deploy skill-free student $\pi_S$. Teacher and
    privilege $z$ are not needed.
\end{enumerate}

\section{What the Guarantee Says}

No formal guarantee is provided. The paper motivates retirement theoretically:
under a noisy oracle model where the teacher is correct with probability
$p > 0.5$ on each token, continuing OPD past the saturation point has
\emph{negative} expected alignment between the OPD gradient and the true
reward gradient. Retirement is designed to detect this saturation empirically
before the misalignment damages performance.

\section{Experimental Findings}

Evaluated on Qwen2.5 at three scales (1.5B, 3B, 7B):

\begin{center}
\small
\begin{tabular}{lcc}
\toprule
Method & ALFWorld & WebShop \\
\midrule
SFT & 54.2 & 59.3 \\
GRPO only & 67.1 & 68.4 \\
Continuous OPD+RL & 71.3 & 72.1 \\
SAGE-OPD & 73.8 & 74.5 \\
\textbf{RetireOPD} & \textbf{85.9} & \textbf{87.4} \\
\bottomrule
\end{tabular}
\end{center}

Success rate (\%), 7B scale. RetireOPD exceeds SAGE-OPD by
$12.1\%$ on ALFWorld and $12.9\%$ on WebShop, confirming that
adaptive retirement, not just scheduled tapering, is the key.

\section{Ablations and Interpretation}

\textbf{No teacher pre-training:} $-8.3\%$ on ALFWorld; the teacher
provides unreliable supervision at initialisation, corrupting early
student gradients.

\textbf{No retirement (continuous OPD):} $-14.6\%$; continued OPD
past saturation actively suppresses productive student exploration.

\textbf{Discrepancy-only retirement:} Premature retirement at 1.5B
scale ($-4.1\%$); the student retires before reaching competence.

\textbf{Success-gap-only retirement:} Over-extended distillation at
7B scale ($-2.7\%$); threshold is met late after OPD stops improving.

\textbf{$\alpha$ sensitivity:} $\alpha \in [0.80, 0.90]$ is robust;
$\alpha < 0.70$ causes premature retirement; $\alpha > 0.95$ delays
retirement past the benefit window.

\section*{Reference}

Yan Yu, Zhengxi Lu, Yizhou Liu, Yichen Pan, Aozhe Wang, Qipeng Chen,
Hua Yang, Wenqi Zhang, Weiming Lu, Qianglong Chen, Yongliang Shen.
\textbf{RetireOPD: Self-Retiring On-Policy Distillation for Agentic
Reinforcement Learning.}
arXiv:2609.20784, September 2026.
\url{https://arxiv.org/abs/2609.20784}

\end{document}
